\section{Exact ball Lam\'e Dirichlet-to-Neumann boundary subsystem}
\label{sec:boundary-gauge-dn}

This section retains the original reference lapse with boundary principal
$f=1$ and restoring parameter $\rho=0$, and imposes the proposed additional
$u_{\scri}=0$ condition. It proves its necessary homogeneous boundary
subsystem. It does not, by itself, prove the interior lapse-wave
initial-boundary value problem or its compatibility with all geometric pole
conditions.

Write $l=1/a$ and $A_{\scri}=S/a$. For the plus-Lie pure physical gauge,
let $X=\tau\partial_t+\xi$, $\eta=\xi+B\tau$, and
\[
 \sigma=\frac{\eta\cdot G+Aw\tau}\Omega,
 \qquad u=A D_0\tau-(A+\Omega)\sigma,
 \qquad v=D_0\eta-l\eta-A^2\nabla\tau.
\]
Regularity of $\xi$ at the fixed sphere gives
$\eta_n=-A_{\scri}\tau$ there. Direct differentiation of the numerator of
$\sigma$ gives
\[
 \sigma_{\scri}=\partial_n\eta_n+l\tau+A_{\scri}\partial_n\tau.
\]
Here $\partial_nG_n=-1/(aS)$ and
$\partial_n(Aw)=-2/a^2$ at $r=S$; both terms are needed.
Consequently $u_{\scri}=0$ implies
$\partial_t\tau=\partial_n\eta_n+l\tau$ and hence
\begin{align}
 \partial_t\eta_n&=-A_{\scri}\partial_n\eta_n+l\eta_n,\\
 \partial_t\eta_A&=-A_{\scri}\partial_n\eta_A+l\eta_A
                         -A_{\scri}\nabla_A\eta_n.
 \label{eq:ball-boundary-eta}
\end{align}
The second equation also uses $v_A=0$ and
$\nabla_A\tau=-\nabla_A\eta_n/A_{\scri}$.

These conditions do not automatically imply the other pole compatibilities:
they give
\[
 \mu_{\scri}=-2A_{\scri}\partial_n\tau,
 \qquad v_{n,\scri}=-A_{\scri}^2\partial_n\tau.
\]
If the original lapse/geometric pole equations additionally require
$\mu_{\scri}=v_{n,\scri}=0$, the further condition is
$\partial_n\tau=0$. Its consistency with the interior time-wave equation
is a separate question. The result below establishes only the
Dirichlet-to-Neumann evolution \eqref{eq:ball-boundary-eta}.

\paragraph{Regular Lam\'e extension.}
Let $\eta$ solve
\[
 \mathcal L\eta=\Delta\eta+\tfrac13\nabla\nabla\cdot\eta=0
 \quad\hbox{in }r<S,
\]
with regularity at the center. Normalize $Y_{\ell m}$ in $L^2(\mathbb S^2)$,
put $L_\ell=\ell(\ell+1)$, and use the poloidal ansatz
\[
 \eta=U(r)Y_{\ell m}n+V(r)\nabla_{\mathbb S^2}Y_{\ell m}.
\]
Its divergence coefficient is
$U'+2U/r-L_\ell V/r$. The vector Laplacian coefficients are
\begin{align}
 (\Delta\eta)_n&=U''+2U'/r-(L_\ell+2)U/r^2+2L_\ell V/r^2,\\
 (\Delta\eta)_A&=V''+2V'/r-L_\ell V/r^2+2U/r^2.
\end{align}
For $(U,V)=r^k(U_0,V_0)$, the divergence is
$r^{k-1}[(k+2)U_0-L_\ell V_0]$. Inserting its gradient yields the two
regular branches, for $\ell\ge1$,
\[
 r^{\ell-1}(U_0,V_0)=r^{\ell-1}(\ell,1),\qquad
 r^{\ell+1}(U_0,V_0)=
 r^{\ell+1}((\ell+1)(\ell-6),\ell+9).
\]
The determinant of their two boundary vectors is $2(7\ell+3)$, so both
Dirichlet components determine a unique regular extension. The factor
$\ell-6$ causes no degeneracy of the extension.

Use normalized tangential amplitude $b=\sqrt{L_\ell}\,V(S)$, so that the
boundary $L^2$ norm is proportional to $|U(S)|^2+|b|^2$.
The ordinary normal derivative has matrix
\[
 \mathcal D_\ell=\frac1{S(7\ell+3)}
 \begin{pmatrix}
 6\ell^2+\ell+3&\sqrt{L_\ell}(\ell-6)\\
 -\sqrt{L_\ell}(\ell+9)&8\ell^2+5\ell-3
 \end{pmatrix}.
\]
In the boundary system the tangential component also contains
$\nabla_A\eta_n$. Adding $\sqrt{L_\ell}/S$ to the lower-left entry gives
\begin{equation}
 \mathcal B_\ell=\frac1{S(7\ell+3)}
 \begin{pmatrix}
 6\ell^2+\ell+3&\sqrt{L_\ell}(\ell-6)\\
 6\sqrt{L_\ell}(\ell-1)&8\ell^2+5\ell-3
 \end{pmatrix}.
 \label{eq:ball-poloidal-B}
\end{equation}
Thus $\partial_t(U,b)^T=G_\ell(U,b)^T$, with
$G_\ell=a^{-1}I-A_{\scri}\mathcal B_\ell$.

\paragraph{Spectrum, conditioning, and low modes.}
The trace is $2\ell/S$ and the eigenvalues of $\mathcal B_\ell$ are
\[
 \frac1S(\ell\pm\delta_\ell),\qquad
 \delta_\ell^2=
 \frac{(\ell-1)(\ell^2-4\ell-3)}{7\ell+3}.
\]
For $\ell=2,3,4$, $\delta_\ell$ is purely imaginary and the generator
has strictly negative real part $-(\ell-1)/a$. For $\ell=5$,
$\delta_\ell^2=4/19$; for $\ell=6$, $\delta_\ell^2=1$.
They too have strictly negative generator eigenvalues.

For every $\ell\ge7$ the uniformly positive metric
\[
 H_\ell=\operatorname{diag}
       \left(\frac{6(\ell-1)}{\ell-6},1\right),
 \qquad I\le H_\ell\le36I,
\]
symmetrizes $\mathcal B_\ell$. Moreover
$0\le\delta_\ell\le(\ell-1)/\sqrt7$. Therefore
\begin{equation}
 \|e^{tG_\ell}\|\le
 6\exp\left[-(1-1/\sqrt7)(\ell-1)t/a\right]
 \qquad(\ell\ge7).
 \label{eq:ball-high-modal-energy}
\end{equation}
This uniform estimate excludes a hidden high-frequency condition-number or
Jordan growth. The finitely many diagonalizable stable modes $2\le\ell\le6$
can be absorbed in a fixed constant; all $\ell\ge2$ consequently admit a
uniform $Ce^{-t/a}$ bound.

For $\ell=1$,
\[
 \mathcal B_1=\frac1S
 \begin{pmatrix}1&-1/\sqrt2\\0&1\end{pmatrix},\qquad
 G_1=\frac1a\begin{pmatrix}0&1/\sqrt2\\0&0\end{pmatrix}.
\]
This is a genuine neutral Jordan block, but it is a fixed finite-dimensional
one: $\|e^{tG_1}\|\le1+t/(\sqrt2a)$. It causes no frequency-dependent
ill-posedness. The $\ell=0$ radial regular solution is $\eta=C r n$ and has
zero boundary generator.

For axial modes $W(r)n\times\nabla_{\mathbb S^2}Y_{\ell m}$, the divergence
vanishes and the regular branch is $W\propto r^\ell$. Their generator is
$-(\ell-1)/a$, with neutral $\ell=1$ rotations and decay for $\ell\ge2$.

\paragraph{Boundary semigroup.}
The normalized vector spherical harmonics are orthogonal. The uniform modal
bound, together with the explicit finite-dimensional neutral blocks, defines
a strongly continuous semigroup on every $H^s$ boundary space and gives
\[
 \|\eta(t)\|_{H^s}\le C_s(1+t/a)\|\eta(0)\|_{H^s}.
\]
On the $\ell\ge2$ subspace it decays uniformly exponentially. The high-mode
generator has size comparable to $\ell$ with a uniformly invertible negative
spectrum, so its domain is $H^{s+1}$ and its stable high-mode part has the
usual order-one analytic smoothing estimates. These claims concern the
explicit boundary operator \eqref{eq:ball-poloidal-B}, not an inferred
energy from the principal symbol alone.

The result is a genuine global-ball Dirichlet-to-Neumann boundary estimate.
It establishes that this necessary repaired homogeneous trace subsystem has
no high-frequency growth. Establishing the interior time-wave problem,
the additional $\partial_n\tau=0$ compatibility when required by the pole
equations, and physical/constraint inhomogeneous forcing remains necessary
before claiming a full gauge or geometric initial-boundary value theorem.
